\documentclass[12pt,a4paper]{report}
\usepackage[utf8]{inputenc}
\usepackage[english]{babel}
\usepackage{amsmath}
\usepackage{amsfonts}
\usepackage{amssymb}
\usepackage{graphicx}
\usepackage{hyperref}
\usepackage{booktabs}
\usepackage{float}
\usepackage{geometry}
\usepackage{setspace}
\usepackage{caption}
\usepackage{subcaption}
\usepackage{titlesec}
\usepackage[T1]{fontenc}
\usepackage[scaled=0.92]{helvet} 
\renewcommand{\familydefault}{\sfdefault}

% Aligning rigorously with the ENSSEA standard
\geometry{margin=2.5cm, left=3cm} 
\setstretch{1.5}

\titleformat{\chapter}[display]
  {\normalfont\huge\bfseries\centering}{\chaptertitlename\ \thechapter}{20pt}{\Huge}
\titleformat{\section}
  {\normalfont\Large\bfseries}{\thesection \ - }{1em}{}
\titleformat{\subsection}
  {\normalfont\large\bfseries}{\thesubsection \ - }{1em}{}
\titleformat{\subsubsection}
  {\normalfont\normalsize\bfseries\underline}{\thesubsubsection \ - }{1em}{}

\begin{document}

\setcounter{chapter}{1} % Set chapter to 2
\chapter{RELATIONSHIPS BETWEEN VARIABLES, LITERATURE REVIEW AND INTERNSHIP FRAMEWORK}

\section*{Introduction of Chapter II}
\addcontentsline{toc}{section}{Introduction of Chapter II}
The transition from generalized, archaic marketing paradigms towards highly targeted, algorithmic enterprise intelligence inherently demands a rigorous theoretical foundation. Prior to deploying empirical machine learning infrastructures—as will be demonstrated profoundly within the structural data of Chapter III—it is fundamentally imperative to establish the strictly defined epistemological boundaries linking the core components of the study. Algorithmic Customer Relationship Management (CRM) does not function as an isolated set of disjointed calculations; it operates as a continuous, interlocking matrix where the outputs of one psychological boundary directly fuel the probabilities of another.

This chapter is entirely devoted to establishing the conceptual framework dictating the profound interaction existing between the study’s primary independent and dependent variables. Specifically, Section 2.1 will meticulously dissect the theoretical literature governing the causal relationship looping Customer Segmentation (RFM \& K-Means) directly into the probabilistic generation of Customer Lifetime Value (CLV). Furthermore, it bridges the micro-economic behavior of individual retail actors into the macro-economic reality of aggregate Sales Forecasting (SARIMA, Prophet, XGBoost). Concluding this section, a synthesized conceptual model is constructed, formulating the exact five research hypotheses (H1–H5) that the subsequent empirical pipeline will mathematically either validate or definitively reject in the context of the Algerian retail sector.

% -----------------------------------------------------------------------------
\section{Relationships between the Variables of the Study}

\subsection{Relationship between Customer Segmentation and Customer Lifetime Value (CLV)}
Customer Lifetime Value (CLV)—defined traditionally within marketing literature as the discounted sum of all future net cash flows generated by an individual consumer—represents the absolute apex of CRM mechanics\footnote{Blattberg, R. C., Getz, G., \& Thomas, J. S. : \textit{Customer Equity: Building and Managing Relationships As Valuable Assets}, Harvard Business School Press, Boston, 2001, p. 88.}. However, mathematical probability models calculating forward-looking yield are notoriously unstable and entirely inaccurate when forcefully applied across a heterogeneous, globally generalized audience. 

The structural relationship between Segmentation and CLV is fundamentally chronological and causal. Statistical probability fields, such as the Beta-Geometric Negative Binomial Distribution (BG/NBD) governing churn, demand behavioral homogeneity internally to prevent outlier vectors from completely rupturing the Gamma distribution parameters\footnote{Fader, P. S., \& Hardie, B. G. S. : \textit{Probability Models for Customer-Base Analysis}, Journal of Interactive Marketing, 23(1), 2009, p. 61.}. Consequently, preliminary heuristic segmentation—specifically the extraction of Recency, Frequency, and Monetary (RFM) components alongside unsupervised K-Means Euclidean partitions—acts as the mandatory precursor variable.

\begin{itemize}
    \item \textbf{Isolation of the Actionable Core}: Continuous tracking studies assert that practically $60\%$ to $70\%$ of a physical retail database constitutes "One-Time" or strictly opportunistic actors. Attempting to formulate probabilistic lifespan tracking on single-purchase actors creates catastrophic arithmetic denominators. Rigorous RFM Segmentation explicitly isolates "Champions" and "Loyal Customers," fundamentally validating the sub-cohorts capable of executing the fractional sub-variables required for Gamma-Gamma average margin tracking.
    \item \textbf{Derivation of the Churn Horizon}: The mathematical definition of a "lost" customer diverges wildly depending on their cluster. A localized "Champion" suddenly missing for 60 days inherently spikes the probability of defection ($P(\text{Alive})$ dropping significantly), whereas a "Dormant" buyer vanishing for 60 days strictly aligns with their historical pattern constraint. Without spatial clustering variables executed concurrently, algorithmic CLV equations inherently remain blind to varying psychological tolerance limits.
\end{itemize}
In summation, Customer Segmentation is the theoretical independent variable that mathematically stabilizes the dependent execution of Customer Lifetime Value parameters.

\subsection{Relationship between CLV and Sales Forecasting: Analytical Interdependencies}
While CLV models evaluate the micro-level lifespan and fractional value of individual human behavioral actors, Sales Forecasting models (such as SARIMA and XGBoost) calculate the macro-level volumetric inventory requirements functioning universally across physical enterprise logistics\footnote{Makridakis, S., Wheelwright, S. C., \& Hyndman, R. J. : \textit{Forecasting: Methods and Applications}, John Wiley \& Sons, New York, 1998, p. 24.}. Within archaic corporate operations, these two dimensions are notoriously structured entirely separate from each other, housed within distinct departmental silos (Marketing versus Logistical Supply Chain).

Modern theoretical frameworks posit that the relationship bridging CLV intimately against aggregate Forecasting is fundamentally interdependent. Macro-level prediction tools dictate \textit{what} logistical volume is expected and \textit{when} the violent seasonal spikes will physically occur. Conversely, the CLV metric identifies precisely \textit{who} is structurally generating that base volume capacity. 

The academic correlation tracking between the variables operates effectively as an acquisition damper mechanism:
\begin{enumerate}
    \item \textbf{Baseline Stabilization via High CLV}: A retail network possessing localized clusters harboring high continuous Median Lifetime Values fundamentally exhibits superior stationary demand curves. High loyalist volume mathematically suppresses extreme time-series variance ($ \epsilon_t $), generating deeply stabilized autoregressive mapping frameworks which allow algorithms like SARIMA to effortlessly anticipate continuous structural baseline cash flows.
    \item \textbf{Forecasting as the Operational Damper}: While probabilistic CLV identifies customers fundamentally maintaining high structural Net Present Value (NPV), this capacity cannot physically manifest functionally if logistical operations incur an "Out-of-Stock" crisis. Precise aggregated temporal forecasting mathematically ensures targeted "Champions" arriving at physical establishments during predicted seasonal weekend spikes natively find inventory completely staged and ready. A failure within the Forecasting variable generates structural supply chain collapse, directly destroying continuous CLV generation execution natively.
\end{enumerate}
Thus, advanced logistics literature verifies that predicting long-term customer relationships requires simultaneous inventory time-series assurance to complete the systemic transactional cycle.

\subsection{Integrated Conceptual Model and Formulation of Research Hypotheses (H1-H5)}
Drawing profoundly upon the intricate mathematical interactions established structurally between heuristic behavior limits, unsupervised spatial boundaries, probabilistic tracking, and macro-aggregate time prediction, an integrated conceptual model natively targeting modern Algerian retail environments can logically be synthesized.

This conceptual framework fundamentally shifts logistical operations entirely away from instinctual estimation, anchoring the core of the PUMA commercial enterprise structure upon strict algorithmic probability execution. To critically validate this specific conceptual integration empirically via Python 3 mechanics in the subsequent Chapter III, five rigorous hypotheses are formally established:

\begin{itemize}
    \item \textbf{Hypothesis 1 (H1)}: Executing strict categorical heuristic RFM rules coupled physically alongside unsupervised continuous log-normalized K-Means geometry isolates specific, statistically predictable behavioral cohorts significantly better than generalized mass audience tracking constraints.
    \item \textbf{Hypothesis 2 (H2)}: Consumer segments algebraically defined via heuristic parameters as "Champions" and "Loyal Customers" inherently display statistically significant higher Median Customer Lifetime Value (CLV) yield projections than cohorts localized within the "Promising" bounds.
    \item \textbf{Hypothesis 3 (H3)}: Operating exclusively within physically brutal, highly turbulent Algerian retail environments defined strictly by violent weekend oscillations, Autoregressive matrices mathematically remembering past sequential errors (SARIMA) comprehensively outperform linear, additive trend frameworks (Prophet) in decreasing continuous absolute volumetric errors natively.
    \item \textbf{Hypothesis 4 (H4)}: Implementing Beta-Geometric decay methodologies specifically accurately flags transitioning "At Risk" actors, isolating precise $P(\text{Alive})$ failure thresholds mathematically prior to totally invisible defection materialization.
    \item \textbf{Hypothesis 5 (H5)}: Constructing an algorithmic, unified interactive Executive Dashboard mathematically overlapping aggregate Sales Forecasting limits intimately against individual CLV parameters structurally generates a fundamentally robust enterprise framework capable of drastically lowering operational Logistical overheads.
\end{itemize}

% -----------------------------------------------------------------------------
\section{Literature Review}

\subsection{Foundational Works on Data-driven Segmentation and ML Clustering}
The theoretical origins of Customer Relationship Management (CRM) relied almost exclusively upon arbitrary demographic divisions (e.g., age, geographic location). However, contemporary marketing literature emphatically declares that demographic variables possess extraordinarily weak correlations concerning future purchasing capacity. True behavioral intelligence emerged heavily in the 1990s through the formalized engineering of the Recency, Frequency, and Monetary (RFM) model. Foundational studies established explicitly that "Recency" remains the single most critically dominant predictor guiding subsequent retail conversion rates\footnote{Hughes, A. M. : \textit{Strategic Database Marketing}, McGraw-Hill, New York, 1994, p. 52.}. Conversely, "Frequency" stabilizes tracking variance, while "Monetary" parameters inherently dictate the mathematical bounds limiting explicit promotional discounting.

However, standard heuristic RFM binning natively generates highly fragmented, disconnected cubes (commonly 125 fractional states) functionally devastating computational efficiency. To combat fractional dispersion, data scientists pivoted aggressively toward embedding RFM continuous parameters into unsupervised Machine Learning (ML) geometries, predominantly anchored by MacQueen’s (1967) \textbf{K-Means Clustering} vector logic\footnote{MacQueen, J. : \textit{Some Methods for Classification and Analysis of Multivariate Observations}, Berkeley Symposium on Mathematical Statistics and Probability, 1967, p. 281.}. 
Modern literature highlights that K-Means effectively minimizes the spatial variance within internal clusters utilizing strictly the sum of squared Euclidean deviations. Empirical modern retail studies dictate that to prevent catastrophic right-skew anomaly warping caused natively by hyper-purchasing "Whale" buyers, spatial dimensions require violent algorithmic standardization executing explicit logarithmic fractional transforms ($\log(1+x)$).

Conversely, spatial tracking literature specifically highlights \textbf{DBSCAN} (Density-Based Spatial Clustering of Applications with Noise) as a historically formidable secondary clustering matrix. Engineered initially as a method mapping abstract spatial shapes heavily fortified against extreme noise\footnote{Ester, M., Kriegel, H. P., Sander, J., \& Xu, X. : \textit{A Density-Based Algorithm for Discovering Clusters in Large Spatial Databases with Noise}, KDD '96 Proceedings, 1996, p. 226.}, contemporary academic reviews strictly demonstrate DBSCAN’s extreme fragility explicitly when predicting retail environments possessing continuous, massively wildly varying global densities. Due inherently to its dependency upon a rigid, static epsilon ($\epsilon$) scanning radius, DBSCAN categorically struggles partitioning commerce data smoothly, commonly shattering the matrix into a massive central glob alongside thousands of mathematically rejected "noise" points, rendering it overwhelmingly inferior to normalized K-Means bounding inside sales architecture.

\subsection{Empirical Studies on CLV Prediction in the Retail Sector}
Historically predicting explicit probability logic monitoring Customer Lifetime Value (CLV) originated as highly fractional contractual modeling—assuming consumers operated identically to subscribed telecommunication users generating rigid sequential income matrices. However, physical sportswear logistics operate natively inside a \textit{non-contractual} sphere; consumers abruptly defect entirely without sequentially modifying central databases.

To resolve strictly non-contractual modeling, Schmittlein, Morrison, and Colombo (1987) engineered the foundational Pareto/NBD framework mathematically mimicking overlapping decay functions. However, the exact mathematical complexity driving the Pareto/NBD model rendered computational calibration operations practically impossible tracking multi-million line matrices.

The absolute disruption mapping non-contractual retail occurred exactly via the seminal publication generating the \textbf{Beta-Geometric/Negative Binomial Distribution (BG/NBD)} by Fader, Hardie, and Lee\footnote{Fader, P.S., Hardie, B.G.S., \& Lee, K.L. : \textit{"Counting Your Customers" the Easy Way: An Alternative to the Pareto/NBD Model}, Marketing Science, 24(2), 2005, p. 275.}. This defining thesis empirically proved that simplifying exact dropout calculations structurally mapping the probability vector utilizing explicit Beta curves alongside isolated Poisson sequences ($P(\text{Alive})$) entirely replicates the immensely robust Pareto/NBD precision—but concurrently demanding absolutely minimal computational RAM architecture.

Subsequently, modern empirical literature demands predicting strictly explicit net margin valuations directly bridging BG/NBD probability via the \textbf{Gamma-Gamma Submodel}. Extensively evaluated mathematically, researchers dictate that calculating fractional Lifetime Value (CLV) relies strictly on ensuring complete statistical independence inherently separating continuous expected physical frequencies ($x$) entirely from their mean physical cash margins ($m$). Assuming Spearman’s Rank tracking empirically calculates statistical coefficients beneath $0.30$, incorporating Gamma distribution yields statistically stabilizes out-of-sample forward-looking testing evaluations massively. Empirical findings practically isolate exactly that acquiring algorithmic intelligence surrounding future baseline yields establishes absolute unbreachable boundaries capping enterprise-level Customer Acquisition Costs (CAC)\footnote{Venkatesan, R. \& Kumar, V. : \textit{A Customer Lifetime Value Framework for Cross-Selling and Customer Acquisition}, Journal of Marketing Research, 41(4), 2004, p. 106.}.

\subsection{Sales Forecasting Benchmarks: State of the Art and Identified Gaps}
Simultaneous to executing individual deterministic lifetimes, accurately predicting aggregated time-series dimensions dictating sequential volumetric logistics bounds defines the remaining foundational predictive constraint characterizing modern economic retail ecosystems.

Literature comprehensively maps historical forecasting heavily relying on purely linear, mathematical \textbf{Autoregressive Integrated Moving Average (ARIMA)} arrays established by Box \& Jenkins (1970). The statistical evolution towards \textbf{SARIMA} explicitly introduced localized harmonic dimensions integrating directly the specific offset $s$, structurally granting pure autoregressive models the mechanical capability executing memory parameters surrounding explicit seasonal cyclicality\footnote{Box, G. E. P. \& Jenkins, G. M. : \textit{Time Series Analysis: Forecasting and Control}, Holden-Day, San Francisco, 1970, p. 112.}. Academic consensus inherently confirms SARIMA provides practically impenetrable mathematical stability, operating perfectly inside violently erratic micro-volatility assuming rigid explicit localized cyclicity loops define the mathematical terrain organically (e.g., explicit continuous weekend traffic anomalies).

Advancing directly into probabilistic supervised learning, Facebook’s computational logistics teams deployed \textbf{Prophet} (2018) generating an explicit additive meta-framework. Engineered strictly resolving structural gaps defining classical matrices, Prophet heavily leverages local regression changepoints modeling localized trend trajectories scaling smoothly directly around disruptive external marketing holidays or completely absent missing logistical strings\footnote{Taylor, S. J. \& Letham, B. : \textit{Forecasting at Scale}, The American Statistician, 72(1), 2018, p. 38.}.
Conversely, gradient-boosting mathematics engineered strictly by Chen & Guestrin generating \textbf{XGBoost} (2016) transitioned continuous matrices completely into arbitrary decision-tree lag classifications. By evaluating explicit temporal sliding-windows (`lag7`, `rolling14`), XGBoost violently outperforms classic algorithms mathematically learning multidimensional non-linear relations completely natively\footnote{Chen, T. \& Guestrin, C. : \textit{XGBoost: A Scalable Tree Boosting System}, KDD '16, ACM, 2016, p. 785.}.

\textbf{Identified Literature Gaps}: 
Despite the mathematically sophisticated global saturation regarding these algorithms natively, a profound academic gap exists executing specifically within the highly idiosyncratic dimensions of physical North African markets (Algeria). Empirical data proves Prophet violently fractures when processing explicit high-frequency localized retail volatility dictating sharp, extreme discontinuous crashes occurring suddenly every single Sunday following monumental Friday volumetric liquidation parameters. Furthermore, while XGBoost maps extreme variance natively utilizing lags, literature isolates categorically its inherent architectural inability projecting extrapolating variables pushing strictly above or beneath the fractional parameters registered physically within its prior internal training limit distributions. 

Consequently, constructing a heavily localized algorithmic comparison formally mapping classical high-frequency SARIMA bounds specifically colliding directly against predictive gradient boosting trees functions absolutely critically directly closing this observable academic deficit operating within an Algerian commercial context.

% -----------------------------------------------------------------------------
\section{Internship Framework : SARL APPAREL (PUMA Algeria)}

\subsection{Company Presentation: History, Activity and Strategic Positioning}
The empirical manifestations of mathematical models only achieve validity when anchored directly into a concrete organizational reality. The professional immersion characterizing this research was conducted within \textbf{SARL Apparel}, the exclusive official authorized distributor of the global conglomerate \textbf{PUMA} in Algeria. 

Founded with a presence in the national sports sector since 2010, Apparel originally operated as a partner for several major international brands. However, a significant strategic pivot occurred in \textbf{2018}, when the enterprise decided to specialize exclusively in the distribution and brand representation of PUMA on the Algerian territory\footnote{Rapport de Stage: \textit{Découverte du milieu professionnel au sein de SARL Apparel}, IFAG, 2021, p. 5.}. 

Today, SARL Apparel manages the entire value chain of the brand locally, encompassing brand image management, multi-channel distribution (B2B and B2C), and a growing E-commerce presence. The company's strategic positioning is anchored on exactly transcribing PUMA's global "Forever Faster" strategy onto the national national territory, ensuring the guarantee of authentic, high-quality products while aggressively developing its retail footprint through the opening of specialized stores.

\subsection{Competitive Environment and the Algerian Retail Context}
The Algerian footwear and clothing industry represents a mature yet volatile sector that has seen numerous brands emerge and vanish over the last two centuries. In this competitive landscape, SARL Apparel has carved a dominant position by focusing on a "short circuit" strategy, prioritizing direct contact and satisfaction of the end customer.

The market environment is characterized by several key dynamics:
\begin{enumerate}
    \item \textbf{Category Dominance}: Within the PUMA portfolio, the \textit{Footwear} segment represents the absolute engine of growth, accounting for over \textbf{45\% of the company's total turnover}. This concentration highlights the necessity for precise stock management and forecasting within this specific category.
    \item \textbf{Digital Shift}: Since July 2020, the company has accelerated its digital transformation, notably through its presence on the \textit{Madina.dz} platform, reaching a diverse target audience of men, women, and children across the country.
    \item \textbf{Socio-Economic Factors}: Despite a relatively low purchasing power in certain segments, there is a marked increase in monthly spending on quality footwear, driven by a youthful population that is highly sensitive to international brand prestige.
\end{enumerate}

\subsection{Field Problem: Available Data, Analytical Needs and Study Scope}
Despite its structured organization—comprising specialized departments in Finance, IT, Logistics, and Marketing—SARL Apparel faces a common challenge in the modern retail era: the transition from "data collection" to "predictive intelligence."

\textbf{Analytical Needs}: The company’s primary objective, as stated in its corporate missions, is to "respond to customer expectations and satisfy them by exploiting all distribution channels." While the SAP/ERP system successfully logs over 180,000 transaction lines annually, these records have historically remained retrospective. There is an urgent analytical need to anticipate customer behavior rather than merely observing it after the purchase.

\textbf{Périmètre de l'étude (Scope)}:
The scope of this research is specifically designed to bridge this gap. By utilizing the 35,696 unique customer profiles identified in the database, our study aims to:
\begin{itemize}
    \item \textbf{Quantify Loyalty}: Moving beyond simple sales tracking to calculate the latent value of each customer through the BG/NBD and Gamma-Gamma models.
    \item \textbf{Optimize Distribution}: Using SARIMA and XGBoost forecasting to ensure that PUMA's retail and e-commerce channels are never hindered by out-of-stock crises during peak demand.
    \item \textbf{Decision Support}: Providing management with an interactive Executive Dashboard that summarizes these complex mathematical vectors into actionable business directives.
\end{itemize}

\section*{Conclusion of Chapter II}
\addcontentsline{toc}{section}{Conclusion of Chapter II}
Chapter II has successfully established the theoretical and organizational framework required to justify the deployment of advanced analytics within SARL Apparel. By examining the literature, we have seen how segmentation, CLV, and forecasting are not isolated metrics but an integrated ecosystem of enterprise survival. 

The transition from the foundational works of RFM and K-Means towards our specific field problematic at PUMA Algeria confirms that "Data is the new oil" only if it is refined through proper methodology. Having identified the company's strategic positioning and the analytical gaps in its current workflow, we are now equipped to transition into Chapter III, where these theoretical hypotheses will be subjected to the rigor of empirical testing and algorithmic validation.

\end{document}
